\documentclass[12pt]{article}
\usepackage[utf8]{inputenc}
\usepackage[T1]{fontenc}
\usepackage{hyperref}
\title{Montenegrin Language Resources for K-12 Curriculum Development}
\author{Compiled by Claude Sonnet 5.5 (claude-sonnet-5-5) for Andr\'as Kornai}
\date{}
\begin{document}
\maketitle

\section{Introduction}
Montenegrin (ISO 639-3 cnr) is the official language of Montenegro. Its status as a language distinct from Serbo-Croatian varieties is politically contested, and sources reflect this. The searches gave limited results and no current ministry textbook list.

\section{Primary Educational Resources}
A four-volume gymnasium series \emph{Crnogorski jezik} (2019, Zavod za udžbenike i nastavna sredstva, Podgorica; Adnan \'Cirgi\'c with co-authors) was reported in a review in \emph{Lingua Montenegrina}. A 2024 adult course \emph{Govorimo crnogorski A1} (same publisher) is in the COBISS.CG catalogue.
\begin{itemize}
\item COBISS.CG entry: \url{https://plus.cobiss.net/cobiss/cg/en/bib/29975812}
\end{itemize}

\section{Linguistic and Reference Materials}
\begin{itemize}
\item CANU, \emph{Rje\v{c}nik crnogorskog narodnog i knji\v{z}evnog jezika} (Register): \url{https://canupub.me/en/knjiga/rjecnik-crnogorskog-narodnog-i-knjizevnog-jezika-registar}
\item University of Oregon Slavic languages guide, Montenegrin section: \url{https://researchguides.uoregon.edu/slavic-languages/montenegrin}
\item Trovesi article on the 2009 \emph{Pravopis} and \emph{Rje\v{c}nik}: \url{https://iris.uniroma1.it/handle/11573/1385902}
\item Zeitschrift f\"ur Balkanologie article on the orthography: \url{https://zeitschrift-fuer-balkanologie.de/index.php/zfb/article/download/288/316/477}
\end{itemize}

\section{Religious and Community Resources}
None were found in these searches.

\section{Digital and Audio Resources}
None were found in these searches.

\section{Curriculum Development Context}
Montenegrin is a language of instruction with a standardized orthography (2009) and state textbook publisher. The Ministry of Education and the Zavod za udžbenike should be consulted for the current approved textbook list.
\end{document}